%%=============================================================================
%% Methodologie
%%=============================================================================

%% TODO: In dit hoofstuk geef je een korte toelichting over hoe je te werk bent
%% gegaan. Verdeel je onderzoek in grote fasen, en licht in elke fase toe wat
%% de doelstelling was, welke deliverables daar uit gekomen zijn, en welke
%% onderzoeksmethoden je daarbij toegepast hebt. Verantwoord waarom je
%% op deze manier te werk gegaan bent.
%%
%% Voorbeelden van zulke fasen zijn: literatuurstudie, opstellen van een
%% requirements-analyse, opstellen long-list (bij vergelijkende studie),
%% selectie van geschikte tools (bij vergelijkende studie, "short-list"),
%% opzetten testopstelling/PoC, uitvoeren testen en verzamelen
%% van resultaten, analyse van resultaten, ...
%%
%% !!!!! LET OP !!!!!
%%
%% Het is uitdrukkelijk NIET de bedoeling dat je het grootste deel van de corpus
%% van je bachelorproef in dit hoofstuk verwerkt! Dit hoofdstuk is eerder een
%% kort overzicht van je plan van aanpak.
%%
%% Maak voor elke fase (behalve het literatuuronderzoek) een NIEUW HOOFDSTUK aan
%% en geef het een gepaste titel.

\chapter{\IfLanguageName{dutch}{Methodologie}{Methodology}}%
\label{ch:methodologie}

\section{Overzicht en werkwijze}
\label{sec:methodologie-overzicht}

Het onderzoek volgt een iteratieve, experimenteel geori{\"e}nteerde aanpak. Gegeven de verkennende aard van dit project, waarbij een nieuw meetsysteem wordt ontwikkeld zonder bestaande standaardmethoden, is een klassieke waterval-methodologie ongeschikt. In plaats daarvan worden opeenvolgende experimentele fasen uitgevoerd, waarbij bevindingen uit elke fase het ontwerp van de volgende fase informeren.

%----------------------------------------------------------------------

\section{Initi{\"e}le verkenning: diepteschatting via ondercamera}
\label{sec:methodologie-verkenning}

Voorafgaand aan de formele onderzoeksfasen werd een eerste exploratieve test uitgevoerd om de haalbaarheid van het oplossingsdomein af te toetsen. Hiervoor werd een iPhone 13 camera, uitgerust met een externe fisheye-lens, onder de trampoline geplaatst om een onderaanzicht te verkrijgen. Aangezien deze test louter exploratief van aard was, werd er in deze fase nog geen lenskalibratie uitgevoerd om de vervorming van de fisheye-lens te corrigeren.

Ondanks de fisheye-lens kon niet de volledige oppervlakte van de trampoline in beeld worden gebracht. Hoewel het grootste deel zichtbaar was, bleven bepaalde randzones buiten het bereik van de camera. Op de opgenomen beelden werd het \textit{Depth Anything v3}-model toegepast, een state-of-the-art model voor monoculaire diepteschatting. De doelstelling was om te onderzoeken of het model de vervorming van het trampolinedoek tijdens een sprong kon vastleggen als indicatie van de landingspositie.

De resultaten waren echter onvoldoende. Aangezien het trampolinedoek uit dit perspectief semi-transparant is, biedt het nauwelijks textuur of visuele referentie voor een dieptemodel. Bijgevolg kon het model geen betrouwbare dieptekaart genereren. Bovendien vervormde het doek op het dieptepunt van de landing zo sterk, dat het overzicht in het beeld verloren ging. Dit maakte een betrouwbare voorspelling van de horizontale verplaatsing mateloos complex. Figuur~\ref{fig:diepteschatting-verkenning} toont een vergelijking tussen een origineel camerabeeld en de bijbehorende dieptemap. Dit leidde tot een herori{\"e}ntering naar een homografiegebaseerde methode via een schuin zijdelings aanzicht.

\begin{figure}[ht]
    \centering
    \includegraphics[width=\linewidth]{depth_anything_test}
    \caption[Vergelijking originele beelden en dieptemap]{Vergelijking tussen originele beelden en de dieptemap via Depth Anything v3.}
    \label{fig:diepteschatting-verkenning}
\end{figure}

%----------------------------------------------------------------------

% !TEX root = CoeneDylanBP.tex
%%=============================================================================
%% Intrinsieke Camerakalibratie
%%=============================================================================

\section{Intrinsieke camerakalibratie}
\label{sec:kalibratie}

Voordat een homografie kan worden berekend, moet de vervorming van de cameralens worden gecorrigeerd. Dit is een essentiële stap, aangezien een homografietransformatie veronderstelt dat rechte lijnen in de echte wereld ook als perfect rechte lijnen in het beeld verschijnen. In de praktijk introduceert vrijwel elke cameralens echter een zekere mate van optische vervorming, waardoor dit niet langer het geval is. Het doel van deze fase is dan ook om de \textit{intrinsieke cameraparameters} nauwkeurig te bepalen. Hiermee kan elk videoframe worden rechtgetrokken voordat de verdere beeldverwerking en positiebepaling plaatsvinden.

De gebruikte aanpak is gebaseerd op de methode van \textcite{Zhang2000}, die ook door OpenCV als standaard wordt gehanteerd. Hierbij wordt een schaakbordpatroon (\textit{checkerboard}) met bekende afmetingen gefilmd vanuit meerdere verschillende hoeken en afstanden. Omdat men exact weet hoe de geometrie van het schaakbord eruit hoort te zien, kunnen de afwijkingen in het camerabeeld worden gemeten. Het algoritme detecteert automatisch de hoekpunten van het schaakbord in elk frame en berekent op basis van de geobserveerde vervorming de volgende parameters:

\begin{itemize}
    \item De \textbf{camerakalibratiematrix} $K$, die de brandpuntsafstand ($f_x, f_y$) en het optische middelpunt ($c_x, c_y$) bevat:
          \[
              K = \begin{bmatrix}
                  f_x & 0   & c_x \\
                  0   & f_y & c_y \\
                  0   & 0   & 1
              \end{bmatrix}
          \]
    \item De \textbf{distortieco{\"e}ffici{\"e}nten}, die beschrijven hoe de lens het beeld vervormt.
\end{itemize}

Na de definitieve bepaling van deze parameters wordt elk frame van de video ontvormd met behulp van de functie \texttt{cv2.undistort} uit de OpenCV-bibliotheek. De kwaliteit van deze kalibratie wordt gemeten aan de hand van de \textit{reprojection error}. Dit is het gemiddelde verschil in pixels tussen de geobserveerde hoekpunten op de foto en de terug berekende positie van diezelfde punten volgens het model. Een waarde kleiner dan 1 pixel wordt hierbij als aanvaardbaar beschouwd voor nauwkeurige metingen.
%%=============================================================================
%% Homografietransformatie
%%=============================================================================

\section{Homografietransformatie: van pixel naar meter}
\label{sec:homografie}

% TODO: Literatuurstudie over formaat van de trampoline
Nu elk videoframe vrij is van lensvervorming, worden de pixelcoördinaten $(u, v)$ uit het beeldvlak via een basistransformatie geprojecteerd naar fysieke coördinaten $(x, y)$ op het trampolinedoek, uitgedrukt in meters. Aangezien de afmetingen van een wedstrijdtrampoline gestandaardiseerd zijn, kan er gebruik worden gemaakt van een \textit{planaire homografietransformatie}.

Een homografie $H$ is een $3 \times 3$ matrix die de projectieve transformatie beschrijft tussen twee vlakken: het camerabeeld en het trampolinedoek.

De transformatie wordt bepaald door vier corresponderende puntenparen:
\begin{enumerate}
    \item De vier hoekpunten van de trampolinemat worden in het videoframe (pixelcoördinaten) bepaald aan de hand van een \textit{landmark detection} model. Dit elimineert de nood aan manuele kalibratie en maakt de aanpak robuuster voor bewegende camerabeelden.
    \item De re{\"e}le afmetingen van de trampoline (bijv. 4.28\,m $\times$ 2.14\,m) worden als doelcoördinaten opgegeven. % TODO: Concrete afmetingen van de trampoline
\end{enumerate}

Met behulp van \texttt{cv2.getPerspectiveTransform} of \texttt{cv2.findHomography} wordt de matrix $H$ berekend. Vervolgens kan elke pixel $(u, v)$ op het trampolinedoek worden omgezet naar een coördinaat $(x, y)$ via:

\[
\begin{pmatrix} x' \\ y' \\ w \end{pmatrix} = H \cdot \begin{pmatrix} u \\ v \\ 1 \end{pmatrix}, \quad (x, y) = \left(\frac{x'}{w},\, \frac{y'}{w}\right)
\]

De nauwkeurigheid van de homografie wordt ge{\"e}valueerd door meetpunten op bekende posities op de mat te projecteren en te vergelijken met de gecorrigeerde re{\"e}le afstand. De verwachte maximale afwijking bedraagt enkele centimeters, afhankelijk van de cameraresolutie en de nauwkeurigheid van de handmatige hoekselectie.

% TODO: Figuur die de homografie toont + plot van de evaluatie
% !TEX root = CoeneDylanBP.tex
%%=============================================================================
%% Objectdetectie en landingsframe detectie
%%=============================================================================

\section{Object- en landingsframe detectie}
\label{sec:yolo}

Met een gekalibreerde camera en berekende homografiematrix is het doel van deze fase om een \textit{first proof-of-concept} te realiseren voor de meting van horizontale verplaatsing (HD). Hiervoor wordt een objectdetectiemodel ingezet om de springer per frame te lokaliseren en zijn relatieve verplaatsing doorheen de sprong te bepalen. Daarnaast wordt geprobeerd om het landingsframe te detecteren aan de hand van de verticale snelheid van het zwaartepunt van de bounding box.

Als detectiemodel wordt YOLOv11 gebruikt~\autocite{Jocher_Ultralytics_YOLO_2023}. Dit model is geschikt voor dit project vanwege de hoge verwerkingssnelheid en de bewezen nauwkeurigheid bij het herkennen van personen in beweging. De volledige pipeline verloopt als volgt:

\begin{enumerate}
    \item YOLOv11 detecteert de springer in elk frame en geeft een bounding box terug.
    \item Het zwaartepunt (middelpunt) van de bounding box wordt berekend in pixelcoördinaten $(u, v)$.
    \item Doorheen de frames wordt de verticale snelheid van het zwaartepunt berekend.
    \item Het moment van landing wordt geschat door te kijken naar het moment dat deze snelheid plotseling afneemt.
\end{enumerate}

Het gebruik van een bounding box zwaartepunt als referentiepunt om de horizontale verplaatsing te bepalen kent een inherente beperking: de grootte en positie van de bounding box vari{\"e}ren afhankelijk van de houding van de springer (bv.\ gestrekte armen vs.\ een gekromde sprong). Dit kan leiden tot een systematische fout in de geprojecteerde positie. Deze beperking geeft een logische aanleiding om in Fase~\ref{sec:pose} over te stappen naar Pose Estimation.

De gemeten landingszones worden vergeleken met manueel bepaalde referentieposities. De nauwkeurigheid wordt uitgedrukt als een percentage correct geclassificeerde landingszones (accuracy) en als een gemiddelde afwijking in meters. % TODO: Dit is nog niet zeker

%%=============================================================================
%% Vluchttijdschatting
%%=============================================================================

\section{ToF-schatting}
\label{sec:tof}

Naast de horizontale verplaatsing is de vluchttijd (ToF) een tweede kerncomponent van de eindscore. Het doel van deze fase is om de ToF en het exacte moment van contact te bepalen op basis van de framerate van de input video, zonder dat er extra sensoren aan te pas komen.

Omdat de gekozen camera een schuin zijdelings aanzicht registreert, ligt de directe schatting van de absolute $z$-co{\"o}rdinaat (hoogte) niet voor de hand. In plaats daarvan wordt de ToF rechtstreeks opgemeten door middel van \textit{contactdetectie} en de bekende video-framerate.

De ToF-inschatting gebeurt met behulp van de landingsframes die in Fase~\ref{sec:yolo} zijn bepaald op basis van de verticale snelheid. Door het frame-verschil tussen opeenvolgende contactmomenten te delen door de framerate van de video, wordt de tijd tussen landing en afstoot berekend. De uiteindelijke score voor vluchttijd wordt bepaald door de som van de tijd tussen deze contactmomenten te berekenen.

De geschatte ToF-waarden worden vergeleken met een manueel bepaalde referentie, afgeleid uit video via frame-telling, en uitgedrukt in:
\begin{itemize}
    \item \textbf{MAE} (Mean Absolute Error): gemiddelde absolute afwijking in seconden.
    \item \textbf{MAPE} (Mean Absolute Percentage Error): relatieve fout t.o.v.\ de referentiewaarde.
\end{itemize}

De doelstelling is een MAE van minder dan 0.2\,s over een volledige reeks van 10 sprongen.

% !TEX root = CoeneDylanBP.tex
%%=============================================================================
%% Verfijning via skeletdetectie
%%=============================================================================

\section{Pose Estimation}
\label{sec:pose}

Indien de baseline uit Fase~3 een voldoende proof-of-concept oplevert, wordt het zwaartepunt vervangen door een nauwkeuriger referentiepunt via \textit{Pose Estimation}. Het doel is de systematische fout door bounding box-variatie te elimineren.

Om de positie langs de korte zijde van de trampoline te bepalen, wordt louter in het geïdentificeerde \textit{landingsframe} (bepaald in Fase~\ref{sec:tof}) een Pose Estimation model ingezet. Dit detecteert specifieke lichaamspunten (\textit{keypoints}) van de springer. Om de landingslocatie te specificeren wordt specifiek gezocht naar de enkels. Het berekend middelpunt tussen deze twee keypoints dient als een nauwkeurig referentiepunt voor de landing van de voeten. Dit elimineert de foutgevoeligheid veroorzaakt door vari{\"e}rende armhouding of lichaamsrotatie van de springer die zich in de volledige bounding box voordoen.

Twee kandidaatmodellen worden ge{\"e}valueerd:
\begin{itemize}
    \item \textbf{ViTPose}~\autocite{Xu2022}: een Vision Transformer-gebaseerd model met state-of-the-art prestaties op \textit{keypoint detection benchmarks}.
    \item \textbf{MediaPipe Pose}~\autocite{Lugaresi2019}: een lichter, real-time alternatief van Google dat geschikt is voor deployment op beperkte hardware.
\end{itemize}

Dit berekende middelpunt in pixelco{\"o}rdinaten wordt vervolgens met behulp van de eerder vermelde homografiematrix $H$ getransformeerd naar het co{\"o}rdinaat in het stelsel van de trampoline. Op basis van deze transformatie wordt bepaald op welke HD-zone de landing wordt geprojecteerd binnen het stelsel van de trampoline. Zo wordt de horizontale verplaatsing nauwgezet vastgesteld.

Doordat het referentiepunt gefixeerd is en niet varieert met armbeweging, wordt een lagere systematische fout verwacht ten opzichte van de bounding box-zwaartepuntbenadering. De vergelijking van beide benaderingen vormt een deelresultaat van het onderzoek.

De validatie verloopt op identieke wijze als in Fase~\ref{sec:yolo} (accuracy en gemiddelde afwijking in meters), maar met de Pose Estimation pipeline als detector.

